\chapter{A Probability Toolkit}
\label{app:toolkit}

This appendix collects standard results used throughout the book, with references for proofs.

\section{Weak convergence}

A sequence of probability measures $\mu_n$ on a Polish space $E$ converges weakly to $\mu$ if $\int f\dd\mu_n\to\int f\dd\mu$ for every bounded continuous $f$. Prokhorov's theorem~\cite{prokhorov1956} states that a family of probability measures on a Polish space is relatively compact in the weak topology if and only if it is tight: for every $\varepsilon>0$ there is a compact $K$ with $\mu(K)\geq1-\varepsilon$ for all members of the family. The portmanteau theorem gives equivalent formulations through open and closed sets, and Skorokhod's representation theorem~\cite{skorokhod1956} realizes weakly convergent sequences as almost surely convergent random variables on a common probability space. Billingsley~\cite{billingsley1999} proves all three. Tightness arguments are the backbone of hydrodynamic limits: one proves that the laws of the empirical density fields are tight, identifies every limit point as a solution of the limiting equation, and concludes by uniqueness.

\section{Metrics on probability measures}

The total variation distance is $\|\mu-\nu\|_{TV} = \sup_A|\mu(A)-\nu(A)|$. The Hellinger distance is $d_{\mathrm{Hell}}(\mu,\nu)^2 = \frac12\int(\sqrt{\dd\mu/\dd\lambda}-\sqrt{\dd\nu/\dd\lambda})^2\dd\lambda$ for any dominating $\lambda$. The bounded Lipschitz distance $d_{BL}$ takes the supremum of $|\int f\dd\mu-\int f\dd\nu|$ over functions with $\|f\|_\infty+\mathrm{Lip}(f)\leq1$ and metrizes weak convergence. The Wasserstein distance $W_p$ is the infimum of $(\E|X-Y|^p)^{1/p}$ over couplings and metrizes weak convergence plus convergence of $p$th moments. Pinsker's inequality $\|\mu-\nu\|_{TV}\leq\sqrt{\frac12\KL(\mu\|\nu)}$ and the bound $d_{\mathrm{Hell}}^2\leq\|\mu-\nu\|_{TV}\leq\sqrt2\,d_{\mathrm{Hell}}$ relate the information theoretic and metric notions. Pinsker's inequality and the comparison of total variation and Hellinger distance are in Section~2.4 of Tsybakov~\cite{tsybakov2009}. The bounded Lipschitz distance and its metrization of weak convergence are in Section 11.3 of Dudley~\cite{dudley2002}, and the Wasserstein distance is treated in Chapter 6 of Villani~\cite{villani2009}.

\section{Martingales}

A martingale with respect to a filtration $(\mathcal{F}_n)$ is an adapted integrable process with $\E[M_{n+1}\mid\mathcal{F}_n] = M_n$. Doob's convergence theorem states that a martingale bounded in $L^1$ converges almost surely, and in $L^1$ if it is uniformly integrable. Doob's maximal inequality gives $\E\sup_{k\leq n}|M_k|^2\leq4\E M_n^2$. The optional stopping theorem states that $\E M_\tau = \E M_0$ for bounded stopping times, and for unbounded ones under uniform integrability. In continuous time, L\'evy's characterization states that a continuous local martingale $M$ with $\langle M\rangle_t = t$ is a Brownian motion, and the Dambis, Dubins and Schwarz theorem states that every continuous local martingale is a time changed Brownian motion. The discrete time results are proved in Kallenberg~\cite{kallenberg2002} and the continuous time results in Revuz and Yor~\cite{revuzyor}, where the Dambis, Dubins and Schwarz theorem is in Section~1 of Chapter~V.

\section{Large deviations}

A family $\mu_n$ satisfies a large deviation principle with rate $n$ and good rate function $I$ if $I$ has compact level sets and for all Borel $A$
\[
-\inf_{A^\circ}I\leq\liminf_n\frac1n\log\mu_n(A)\leq\limsup_n\frac1n\log\mu_n(A)\leq-\inf_{\bar A}I.
\]
The contraction principle transfers a large deviation principle through continuous maps, with rate function $\tilde I(y) = \inf\{I(x):F(x) = y\}$. Varadhan's lemma computes exponential moments, and the G\"artner Ellis theorem derives a principle from the limit of normalized cumulant generating functions. Cram\'er's, Sanov's~\cite{sanov1957}, Schilder's~\cite{schilder1966} and the Donsker Varadhan~\cite{donsker1975} theorems are the principles for sample means, empirical measures, Brownian paths scaled by $\sqrt\varepsilon$, and occupation measures of Markov processes respectively. Schilder's rate function $\frac12\int_0^1|\dot\phi|^2\dd t$ is the Cameron Martin norm, and its finiteness defines the Cameron Martin space of Chapter~\ref{ch:nsinference}. Dembo and Zeitouni~\cite{dembozeitouni} prove the contraction principle (Section~4.2.1), Varadhan's lemma (Section~4.3) and the G\"artner Ellis theorem (Section~2.3).

\section{Semigroups and generators}

A Markov semigroup $(P_t)$ on bounded functions has generator $\mathcal{L}f = \lim_{t\downarrow0}(P_tf-f)/t$ on its domain, and the Kolmogorov equations read $\partial_tP_tf = \mathcal{L}P_tf = P_t\mathcal{L}f$. A probability measure $\mu$ is invariant if $\int P_tf\dd\mu = \int f\dd\mu$, and reversible if $P_t$ is self adjoint in $L^2(\mu)$. For a reversible semigroup the Dirichlet form $\mathcal{E}(f,f) = -\int f\mathcal{L}f\dd\mu$ controls convergence: a Poincar\'e inequality $\Var_\mu(f)\leq C\,\mathcal{E}(f,f)$ gives $L^2$ convergence at rate $e^{-t/C}$, and a log Sobolev inequality $\Ent_\mu(f^2)\leq2C\,\mathcal{E}(f,f)$ gives entropy convergence at rate $e^{-2t/C}$. The Bakry \'Emery criterion $\Gamma_2\geq K\Gamma$ implies both with $C = 1/K$, where $\Gamma(f) = \frac12\mathcal{L}f^2-f\mathcal{L}f$ and $\Gamma_2(f) = \frac12\mathcal{L}\Gamma(f)-\Gamma(f,\mathcal{L}f)$; for the diffusion with generator $\Delta-\nabla V\cdot\nabla$ on a Riemannian manifold, $\Gamma_2\geq K\Gamma$ is the condition $\Ric+\nabla^2V\geq K$. Ethier and Kurtz~\cite{ethier1986} give the semigroup theory, Bakry, Gentil and Ledoux~\cite{bakry2014} the functional inequalities (Chapters 4 and 5), and the criterion is due to Bakry and \'Emery~\cite{bakry1985}.
